\section{Mail Delivery}
\label{sec:cses-mail-delivery}

\problemstatement{Given an undirected graph of $n$ nodes and $m$ edges,
find a route that starts and ends at node $1$ and traverses \emph{every
edge exactly once} (an Eulerian circuit), or output \texttt{IMPOSSIBLE}.}

\begin{patternbox}
An \textbf{Eulerian circuit} exists iff the graph is connected (over
edge-bearing nodes) and \emph{every vertex has even degree}. When it
exists, \textbf{Hierholzer's algorithm} stitches it together in linear
time.
\end{patternbox}

\subsection*{Degree check, then Hierholzer}

First verify every vertex has even degree and all edges lie in one
connected piece reachable from node $1$; otherwise no circuit exists.
Hierholzer's then walks edges greedily from node $1$, marking each used,
until it returns to a node with no unused edges -- backtracking (popping to
an output list) whenever stuck, which merges sub-cycles into one trail.
Reversing the popped order yields the Eulerian circuit. A per-node pointer
skips already-used edges so total work stays linear.

\begin{ideabox}[Even Degrees Guarantee You Can Always Leave]
Entering a vertex uses one incident edge; an even degree guarantees an
unused edge to leave by, so the walk only gets stuck back at the start.
Hierholzer exploits this to splice every local cycle into a single circuit.
\end{ideabox}

\subsection*{Algorithm}

\begin{enumerate}
  \item Check all degrees even and edges connected; else
        \texttt{IMPOSSIBLE}.
  \item Hierholzer from node $1$ using a stack and per-node edge pointer,
        marking edges used.
  \item Reverse the popped nodes; if it uses all $m$ edges, print it.
\end{enumerate}

\subsection*{Dry run}

\dryrun{Square $1\!-\!2$, $2\!-\!3$, $3\!-\!4$, $4\!-\!1$.}

\begin{itemize}
  \item All degrees $2$ (even), connected -- a circuit exists.
  \item Hierholzer yields $1\to2\to3\to4\to1$.
\end{itemize}

\subsection*{C++ implementation}

\begin{lstlisting}
#include <bits/stdc++.h>
using namespace std;

int main() {
    int n, m;
    cin >> n >> m;
    vector<vector<pair<int,int>>> adj(n + 1);    // {neighbor, edgeId}
    vector<int> deg(n + 1, 0);
    for (int i = 0; i < m; i++) {
        int a, b; cin >> a >> b;
        adj[a].push_back({b, i});
        adj[b].push_back({a, i});
        deg[a]++; deg[b]++;
    }
    for (int i = 1; i <= n; i++)
        if (deg[i] % 2 != 0) { cout << "IMPOSSIBLE\n"; return 0; }

    vector<int> ptr(n + 1, 0);
    vector<bool> used(m, false);
    vector<int> circuit;
    stack<int> st;
    st.push(1);
    while (!st.empty()) {
        int u = st.top();
        while (ptr[u] < (int)adj[u].size() && used[adj[u][ptr[u]].second]) ptr[u]++;
        if (ptr[u] == (int)adj[u].size()) { circuit.push_back(u); st.pop(); }
        else {
            auto [v, id] = adj[u][ptr[u]];
            used[id] = true;
            st.push(v);
        }
    }
    reverse(circuit.begin(), circuit.end());
    if ((int)circuit.size() != m + 1) { cout << "IMPOSSIBLE\n"; return 0; }
    for (int v : circuit) cout << v << " ";
    cout << "\n";
    return 0;
}
\end{lstlisting}

\begin{complexitybox}
\textbf{Time Complexity.} $O(n+m)$ -- each edge is used once; pointers avoid
rescanning. \textbf{Space Complexity.} $O(n+m)$.
\end{complexitybox}

\begin{takeawaybox}
Eulerian circuits are a degree check plus Hierholzer's stitching of cycles.
The per-node pointer and edge-used markers keep it linear even with
multi-edges. The circuit-size check at the end catches disconnection.
\end{takeawaybox}
